\subsubsection{The archimedean mean on the reference chain}
\label{subsec:selected-period-means}

The single-entry estimate involves both the CM period and the norm of a theta vector. The squared kernel height uses their means over the same embeddings. Write \(K=\mathbb Q(\sqrt{-d})\) for the CM curve field,
\(r_E\) for the root discriminant of the CM field \(E\), and
\(\ell_K=(L'/L)(1,\chi_{-d})\). The discriminant of \(K\) is
\(-d\), since \(d\equiv7\bmod8\) is prime. In this paragraph an
overline on a real quantity means its average over the embeddings
of a number field defining the curve, its differentials, the
finite class data and the coefficients of the selected base
sections.

The field of definition can be enlarged to contain the torsion points used
in a calculation. If \(F'/F\) is such an enlargement and \(a_\rho\) is a quantity
attached to an embedding of \(F\), then
\[
 \frac{1}{[F':\Q]}\sum_{\rho':F'\hookrightarrow\C}a_{\rho'|_F}
    =\frac{[F':F]}{[F':\Q]}\sum_{\rho:F\hookrightarrow\C}a_\rho .
\] In particular one may first average
above the infinite place of the period matrix. Every embedding
there bounds the same squared operator norm, and the two
conjugate embeddings of its totally imaginary coefficient field
give that place with the same multiplicity.

\begin{proposition}[Mean normalized differential area]\label{prop:mean-differential-area}
For the full-order CM models of \(K=\mathbb Q(\sqrt{-d})\)
with discriminant-one period lattice as in
Proposition~\ref{prop:level-free-theta-jets}, the mean normalized
area of the axis
generator \(d_\tau\) is
\[
 \overline{\log\frac{A_\rho}{\pi|\rho(d_\tau)|^2}}
   =\gamma+\log(2\pi/r_E)-\ell_K .
\]
The mean is unchanged by adjoining the finite class or
coefficient data of the selected test.
\end{proposition}

\begin{proof}
\proofstep{Average the discriminant-one area over the full ideal class group.}
Let \(A_\rho\) be the area of the period lattice at embedding
\(\rho\), with its discriminant normalized to one. The main
theorem of complex multiplication for the full order of \(K\)
makes its ideal classes equidistributed over these embeddings.
Indeed over \(K\) the Galois action on the singular moduli is
the ideal class action on the lattices, simply transitive. Complex
conjugation permutes the same classes, and adjoining the
coefficients of a discriminant-one equation gives equal
multiplicities. A homothety \(\Lambda\mapsto c\Lambda\) changes the
discriminant by \(c^{-12}\) and the area by \(|c|^2\). For two
discriminant-one representatives in the same class,
\[
 c^{-12}=1\quad\Longrightarrow\quad
 \operatorname{area}(c\Lambda)=\operatorname{area}(\Lambda).
\] \cite{cm-reciprocity}

\proofstep{Compute the average using the first limit formula.}
Choose a positively oriented basis \(\omega,\omega\tau\) in a
class, write \(y=\Im\tau\), and put
\(Z(\tau,s)=\sum' y^s/|m+n\tau|^{2s}\). Kronecker's first
limit formula gives residue \(\pi\) at one and constant term
\(2\pi(\gamma-\log 2-\log(\sqrt y\,|\eta(\tau)|^2))\).

For an ideal \(I\), its unscaled covolume is
\(\mathrm{N}(I)\sqrt d/2\). The sum over its nonzero elements
counts each integral ideal in the inverse class once for each
unit, by \(a\mapsto(a)I^{-1}\). Hence, for \(\Re s>1\),
the sum of the \(Z\)'s over the classes is
\[
       |\mathcal O_K^\times|(\sqrt d/2)^s\zeta(s)L(s,\chi_{-d}).
\]
This also identifies their Laurent coefficients at one. Divide the
constant term by the residue on both sides. The result is
\[
  2\,\overline{\log(\sqrt y\,|\eta(\tau)|^2)}
       =\gamma-\log(2\sqrt d)-\ell_K .
\]
The analytic discriminant is
\((2\pi/\omega)^{12}\eta(\tau)^{24}=1\), so
\(A_\rho=(2\pi)^2y|\eta(\tau)|^4\). We obtain
\begin{equation}\label{eq:selected-curve-area}
        \overline{\log(A_\rho/\pi)}
             =\gamma+\log(2\pi/\sqrt d)-\ell_K .
\end{equation}
\cite{kronecker-period-mean}\label{cor:selected-period-means-1}

There is a second factor from the host. With
\(\mathfrak D=\mathfrak D_{E/K}\) and the inverse half different
\(I=\mathfrak b^{-1}\), equation~\eqref{eq:host-different-factors}
gives \(\mathfrak D=\mathfrak b\bar{\mathfrak b}\) for the fixed target
character, including the primes of its conductor. Its relative
trace dual for the bilinear pairing is
\(I^*=\mathfrak D^{-1}I^{-1}=\bar{\mathfrak b}^{-1}\). In the split coordinates of
\(E\otimes_K M\), for a field \(M\) containing its embeddings,
the pairing is the sum of coordinate products. Tensor the
perfect pairing of \(I\) and \(I^*\) with \(\mathcal O_M\)
over \(\mathcal O_K\). These are projective modules, so the
pairing stays perfect.

If \(e_\tau\) is the coordinate vector on the \(\tau\)-axis,
then pairing \(ce_\tau\) with \(y\in I^*\) gives \(c\tau(y)\).
Perfectness thus identifies the ideal of allowed values as
\[
 \{c:c\tau(I^*)\subseteq\mathcal O_M\}
      =(\tau(I^*)\mathcal O_M)^{-1}
      =\tau(\bar{\mathfrak b})\mathcal O_M .
\]
It is principal for the defining field of the selected
generators \(d_\tau\). This interpretation persists after
enlargement, and descends to any defining field by faithful flatness.
Since the two conjugate factors have equal absolute ideal norm,
the ideal norm formula and the tower formula for discriminants give
\[
  \overline{\log |\rho(d_\tau)|^2}
   =\frac{\log\mathrm N_{E/\mathbb Q}\mathfrak D}{[E:\mathbb Q]}
   =\log r_E-\log\sqrt d .
\]
Subtracting from \eqref{eq:selected-curve-area} gives the
differential term with its metric and its sign,
\begin{equation}\label{eq:selected-generator-mean}
  \overline{\log\frac{A_\rho}{\pi|\rho(d_\tau)|^2}}
     =\gamma+\log(2\pi/r_E)-\ell_K=\log\Omega .
\end{equation}
This holds at each individual type label. The same calculation
at the enhanced finite primes explains why the axis generator
there has valuation zero.
\end{proof}

\paragraph{The squared norm of the theta vector.}
\begin{proposition}[Mean squared theta norm on the reference chain]\label{prop:mean-reference-theta-norm}
For the discriminant-normalized curve with its marked
nonzero two-torsion point, \(N_\rho\) is the normalized
Gaussian squared norm of its even theta section. Its mean satisfies
\[
 \overline{\log N_\rho}\le\tfrac12(\log\sqrt d-\gamma+\ell_K).
\]
The positive-part mean has the additional allowance
\(\tfrac14\log d+2\log C_0\), with \(C_0\) an absolute bound
for the even theta constants of the reduced ideal bases.
The squared norm on the standard Hom pair at step \(j\)
is \(N_\rho^{D_j}\) over each conjugate model.
\end{proposition}

\begin{proof}
\proofstep{Evaluate the Gaussian integral in the marked characteristic.}

Let \(h_\rho\) be the distinguished nonzero two-torsion point,
and let \(\theta_\rho\) be the even theta function with this
zero, normalized at the origin. In the basis above let
\(\vartheta\) denote the corresponding even Jacobi theta
constant. Its Gaussian squared norm, using normalized Haar
average, is
\[
      N_\rho=(2y)^{-1/2}|\vartheta(0,\tau)|^{-2}.
\]
Indeed integrate the squares of its Fourier series terms over
one real period. Distinct frequencies are orthogonal. For the characteristic shift \(a\in\{0,1/2\}\), integration
in the imaginary period leaves
\[
 \sum_{n\in\mathbb Z}\int_0^1e^{-2\pi y(n+a+t)^2}\,dt
    =\int_{\mathbb R}e^{-2\pi y u^2}\,du=(2y)^{-1/2}.
\] Normalization at the
origin divides by the squared theta constant.

\proofstep{Use the algebraic integer from the half-period.}

The quotient \(|\vartheta/\eta|^{12}\) is the absolute value
of a conjugate of a single nonzero algebraic integer. To verify
the assertion with this half-period, put \(e=\wp(h)\) and
\(Y=(12e^2-g_2)/4\). For the other roots \(e',e''\) of the Weierstrass cubic,
differentiating \(4(X-e)(X-e')(X-e'')=4X^3-g_2X-g_3\) at \(e\)
gives
\[
 (e-e')(e-e'')=(12e^2-g_2)/4=Y .
\] The three half-period difference formulas and
\(\theta_2\theta_3\theta_4=2\eta^3\), together with
discriminant one, give
\(
        |\rho(Y^{-3})|=|\vartheta(0,\tau)/\eta(\tau)|^{12}.
\)
For instance at \(h=\omega/2\) the two differences are
\((\pi/\omega)^2\theta_3^4,(\pi/\omega)^2\theta_4^4\)
and the characteristic is that of \(\theta_2\).

The permutations
of the roots give the other two cases with signs. In the
identity \(u=-Y^{-3}=s_i(\theta_i/\eta)^{12}\), the multipliers
are \((s_2,s_3,s_4)=(-1,1,-1)\). These are multipliers of
the complex ratios at the chosen half-periods. Each conjugate satisfies the monic integral equation over
\(\mathbb Z[j]\)
\[
 (u-16)^3-ju=u^3-48u^2+(768-j)u-4096=0 .
\] This follows on substituting
\(\lambda=\theta_2^4/\theta_3^4\) in
\(j=256(1-\lambda+\lambda^2)^3/(\lambda^2(1-\lambda)^2)\).
Thus \(Y^{-3}\) is integral since \(j\) is an algebraic integer.
Both \(e\) and the discriminant-one model conjugate with the
specified two-torsion datum. The product formula gives
\(\overline{\log|\vartheta|}\ge\overline{\log|\eta|}\).

Combining the Gaussian identity with the area formula yields
\[
        \overline{\log N_\rho}
        \le\tfrac12(\log\sqrt d-\gamma+\ell_K).
\]
A reduced ideal basis gives \(1/2\le y\le\sqrt d/2\)
by the reduced binary quadratic form of discriminant \(-d\).
The defining Gaussian series bounds each even theta constant
at zero by an absolute \(C_0\ge1\) on this range. Thus the negative part obeys the pointwise estimate
\[
 (-\log N_\rho)_+
 =\big(\tfrac12\log(2y)+2\log|\vartheta(0,\tau)|\big)_+
 \le \tfrac14\log d+2\log C_0 .
\]
Adding its mean to \(\overline{\log N_\rho}\) bounds the mean
of \((\log N_\rho)_+\). This permits a coefficient whose
sign can vary elsewhere in the height comparison.
\cite{elliptic-constants}\label{cor:selected-period-means-2}

\proofstep{Transport the normalized one-variable section.}

For the standard source and target the rank of the quadratic
CM Hom space is \(D_j=f N_{j-1}N_j\). We spell out the use of
the theta normalization under conjugation. On its fourfold sum
the normalized vector of Section~\ref{subsec:theta-normalization}
is a normalized finite sum of Heisenberg translates of the
product of \(4D_j\) one-variable sections. At the conjugate
curve the one-variable factor is again the even section
normalized at the origin with zero \(h_\rho\). Let \(s\) denote this normalized section in its local coordinate. Evenness and the logarithmic derivative formula give
\[
 s'(0)=0,\qquad (\log s)''(z)=-\wp(z+h_\rho)-s_2 .
\]
Formula~\eqref{eq:cm-isogeny} computes \(s_2\) algebraically from the nonzero points of a CM isogeny with \(a^2\ne a\bar a\). These identities conjugate
coefficientwise and determine the entire Taylor series.

\proofstep{Compare translated sections and the fourth tensor.}

Write \(s\) for the unshifted section and \(s_t\) for its
translate by a point \(t\) of order \(m\). The quotient has rational
\(m\)-th power, and
\[
            \operatorname{div}\bigl((s_t/s)^m\bigr)
                      =m([h-t]-[h]).
\] After conjugation it is a scalar times the
corresponding rational function on the conjugate curve, by
equality of divisors. More precisely the functions can be taken
over the algebraic coefficients. Their local Laurent series is
algebraic coefficientwise, and the space of rational functions
with that divisor is a line defined over the field of the torsion
points. The product of this rational quotient over the \(m\)
translations by \(jt\) is a root of unity by the period and
commutator laws, as in Section~\ref{subsec:baseline-jets}.

If \(f_1=c f_2\) are the two rational quotients with the
same divisor after conjugation, their translated products differ by
\[
 \prod_{j=0}^{m-1} f_1(z+jt)
     =c^m\prod_{j=0}^{m-1} f_2(z+jt).
\]
Both products are roots of unity in this normalization, hence
\(|c|^m=1\) at every embedding.
Taking the formal \(m\)-th root gives the translated conjugate
section up to a unit of complex modulus, including a translate
vanishing at zero. On the finite quotient, let \(\kappa(x,a)\) be the commutator
character of the translated summand \(G_a\) under the unitary
translation \(W_x\), with \(a\) the corresponding torsion class.
For \(a\ne b\), perfection of the pairing gives \(x\) such that
\(\kappa(x,a)\ne\kappa(x,b)\). Thus
\[
 \langle G_a,G_b\rangle
  = \kappa(x,a)\overline{\kappa(x,b)}
       \langle G_a,G_b\rangle=0 .
\]
The transported classes remain distinct and their commutator
pairing is the conjugate Weil pairing. The same calculation
therefore gives orthogonality after conjugation.

Changing the
phases of the individual summands preserves the norm. The rational CM isometry conjugates to an
isometry, since its adjoint equations are algebraic over the
CM fields, on which conjugation commutes with the embeddings.
Its change of variables preserves normalized averages.
For the index \(I\) of this theta-line transfer the conjugate
section is, up to the determinant scalar of modulus one in
\eqref{eq:kappa-determinant}, \(I^{-1/2}\sum_{a=1}^I\zeta_a G_a\).
Each \(|\zeta_a|=1\), and the \(G_a\) are orthogonal with
squared norm \(N_\rho^{4D_j}\) since each is a unitary translate
of that one product section. Hence
\[
      \left\| I^{-1/2}\sum_{a=1}^I\zeta_a G_a\right\|^2
               = I^{-1}\sum_{a=1}^I\|G_a\|^2
               = N_\rho^{4D_j}.
\] Taking fourth roots of squared
tensor norms now gives \(N_\rho^{D_j}\) for the squared norm
of the standard vector. The ideal-character multiplier on the
standard pair is one.

\end{proof}

\paragraph{Averaging the selected error from the kernel.}
\begin{proposition}[A uniform bound for the mean logarithmic lift norm]\label{prop:mean-reference-lift-price}
For the selected word and reference chain, write \(q_{jz}=\sum_i\nu_{jzi}\) for the size of the partition at lift \(j\) and real place \(z\). Set \(\varphi(u)=u(\log u-1)\) for \(u>0\) and \(\varphi(0)=0\). The mean logarithmic lift bound \(h_B(j)\) satisfies
\[
 h_B(j)\le
 \sum_z\left(\sum_i\varphi(\nu_{jzi})-q_{jz}\log\Omega\right)
       +C fN_m^2(L_a+m^3L).
\]
The constant is independent of the CM field, word, embedding and lift index \(j\)
for fixed \(C_1\), and the estimate applies at both coupled
indices with the given normalization.
\end{proposition}

\begin{proof}
\proofstep{Start from the squared entry bound.}

Put \(n=N_{j-1}, n'=N_j, v=r_E/2\), \(L_a=\log(a r_E)\),
\(q_{jz}=\sum_{i=1}^n\nu_{jzi}\) and
\(\varphi(s)=s(\log s-1)\), with \(\varphi(0)=0\).
For each embedding inducing a given infinite place of the
coefficient field, the single entry in that place has the
logarithmic squared bound
\begin{align*}
 k_j(\rho)={}&\sum_z\left(\sum_i\varphi(\nu_{jzi})
       -q_{jz}\log\frac{A_\rho}{\pi|\rho(d_z)|^2}\right)
      + D_j\log N_\rho + D_j\log(8+32D_jv)\\
 &\quad +2C_1 U_u+
           \sum_z\bigl(1+2nn'\log(2q_{jz}+2)\bigr).
\end{align*}
Here \(C_1=C+4\) by the pointwise estimate in
Section~\ref{subsec:selected-kernel-jets}.
The ordinary and coupled sums now use the multiplicities just
proved. Equivalently one may put \(C_1+58\) in the last bound
for each lift. The upper bound for each lift also includes the
nonnegative logarithms of its two class and two index counts.
These add at most \(112U_u\) per lift if retained explicitly.

\proofstep{Insert the parameter bounds and the two archimedean means.}

Since \(\log g_i\le9L+1\), \(\log T_u\le25mL\) and
\(L,m\ge1\), we have \(U_u\le45fm^3N_m^2 L\).
Indeed,
\[
 \log(g_0g_1T_u)\le18L+2+25mL\le45mL .
\]
Moreover \(fN_m\le b\), \(q_{jz}\le3aN_m\),
\(4x\le b\), \(a=xb\) and \(r_E\ge1\). These give
\(\log(8+32D_jv)=O(L_a)\) and
\(\log(2q_{jz}+2)=O(L_a)\) with absolute constants.

The selected bounds on \(\log d\) and \(|\ell_K|\) are at
most \(L\). Applying the two archimedean means above shows
that the averaged lift price is at most
\[
  \sum_z\left(\sum_i\varphi(\nu_{jzi})-q_{jz}\log\Omega\right)
       +C fN_m^2(L_a+m^3L).
\]
This is the numerical height contribution \(h_B(j)\) on the
same word and at both coupled indices, with \(C\) independent
of the host, word, embedding, and \(j\), for fixed \(C_1\).
The terminal evaluation follows in the next subsection. The
row scalar enters the comparison with the word height.
\label{cor:selected-period-means-3}
\end{proof}
